% style-guide.tex
% Bitcoin Austria Präsentations-Style-Guide
%
% Bauen: cd .. && make PRESENTATION=style

\documentclass[aspectratio=169,t,9pt]{beamer}

\usetheme[progressbar=foot]{metropolis}
\usepackage{../styles/bitcoin-austria}

\title{}
\author{}
\date{}
\institute{}

\begin{document}

% ── Folie 1: Titelfolie ───────────────────────────────────────────────────────
\begin{frame}[plain]
	\begin{tikzpicture}[remember picture, overlay]
		% Full dark background
		\fill[badark] (current page.south west) rectangle (current page.north east);
		% Red accent strip on left
		\fill[bared] (current page.south west)
		rectangle ([xshift=0.32\paperwidth] current page.north west);
	\end{tikzpicture}
	%
	\begin{columns}[T]
		% Left red column — intentionally empty (brand colour strip)
		\begin{column}{0.32\textwidth}\end{column}
		% Right dark column — logo + title text
		\begin{column}{0.65\textwidth}
			\vspace{0.06\paperheight}
			% Horizontal logo (icon + "Bitcoin Austria" text), right-aligned
			\begin{flushright}
				\includegraphics[height=5ex, keepaspectratio]%
				{../logos/logo-bitcoin-austria-light.pdf}\hspace{2pt}
			\end{flushright}
			\vspace{0.05\paperheight}
			{\BlinkHeadline\fontsize{20}{24}\selectfont\color{balight}
				Bitcoin Austria\\[0.2em]
				Präsentations-Style-Guide\par}
			\vspace{0.6em}
			{\fontsize{10}{13}\selectfont\color{bateal}
				Farbpalette · Typografie · Layout-Makros\par}
			\vspace{1.0em}
			{\fontsize{7.5}{10}\selectfont\color{balight!55!badark}
				\texttt{styles/bitcoin-austria.sty} ·
				Blinker (SIL OFL) · XeLaTeX\par}
		\end{column}
	\end{columns}
\end{frame}

% ── Folie 2: Farbpalette ──────────────────────────────────────────────────────
\begin{frame}{Farbpalette}
	\vspace{0.3em}
	\begin{center}
		% Use plain cm coordinates — do NOT use x= scaling (it also scales ++ offsets)
		\begin{tikzpicture}[x=1cm, y=1cm]
			% Swatch dimensions: 2.0cm wide, 1.3cm tall, 0.28cm gap
			% 4 columns at x = 0, 2.28, 4.56, 6.84
			% 2 rows: row 0 at y=0..1.3,  row 1 at y=1.58..2.88

			% ── Row 1 (top): Primärfarben ──
			\fill[bared,    rounded corners=4pt] (0.00,1.58) rectangle ++(2.0cm,1.3cm);
			\fill[baorange, rounded corners=4pt] (2.28,1.58) rectangle ++(2.0cm,1.3cm);
			\fill[badark,   rounded corners=4pt] (4.56,1.58) rectangle ++(2.0cm,1.3cm);
			\fill[bagray,   rounded corners=4pt] (6.84,1.58) rectangle ++(2.0cm,1.3cm);

			\node[text=white, font=\fontsize{7.5}{9}\selectfont\bfseries] at (1.00,2.53) {\texttt{\textbackslash bared}};
			\node[text=white, font=\fontsize{6.5}{8}\selectfont]          at (1.00,2.22) {\texttt{\#E3000F}};
			\node[text=white, font=\fontsize{7.5}{9}\selectfont\bfseries] at (3.28,2.53) {\texttt{\textbackslash baorange}};
			\node[text=white, font=\fontsize{6.5}{8}\selectfont]          at (3.28,2.22) {\texttt{\#F7931A}};
			\node[text=white, font=\fontsize{7.5}{9}\selectfont\bfseries] at (5.56,2.53) {\texttt{\textbackslash badark}};
			\node[text=white, font=\fontsize{6.5}{8}\selectfont]          at (5.56,2.22) {\texttt{\#222222}};
			\node[text=white, font=\fontsize{7.5}{9}\selectfont\bfseries] at (7.84,2.53) {\texttt{\textbackslash bagray}};
			\node[text=white, font=\fontsize{6.5}{8}\selectfont]          at (7.84,2.22) {\texttt{\#4D4D4D}};

			% ── Row 2 (bottom): Akzentfarben ──
			\fill[balight,  rounded corners=4pt] (0.00,0.00) rectangle ++(2.0cm,1.3cm);
			\draw[bagray!50, rounded corners=4pt, line width=0.4pt] (0.00,0.00) rectangle ++(2.0cm,1.3cm);
			\fill[bateal,   rounded corners=4pt] (2.28,0.00) rectangle ++(2.0cm,1.3cm);
			\fill[bablue,   rounded corners=4pt] (4.56,0.00) rectangle ++(2.0cm,1.3cm);
			\fill[bayellow, rounded corners=4pt] (6.84,0.00) rectangle ++(2.0cm,1.3cm);

			\node[text=badark, font=\fontsize{7.5}{9}\selectfont\bfseries] at (1.00,0.95) {\texttt{\textbackslash balight}};
			\node[text=badark, font=\fontsize{6.5}{8}\selectfont]          at (1.00,0.64) {\texttt{\#FEFEFE}};
			\node[text=white, font=\fontsize{7.5}{9}\selectfont\bfseries]  at (3.28,0.95) {\texttt{\textbackslash bateal}};
			\node[text=white, font=\fontsize{6.5}{8}\selectfont]           at (3.28,0.64) {\texttt{\#73BFB8}};
			\node[text=white, font=\fontsize{7.5}{9}\selectfont\bfseries]  at (5.56,0.95) {\texttt{\textbackslash bablue}};
			\node[text=white, font=\fontsize{6.5}{8}\selectfont]           at (5.56,0.64) {\texttt{\#0D579B}};
			\node[text=badark, font=\fontsize{7.5}{9}\selectfont\bfseries] at (7.84,0.95) {\texttt{\textbackslash bayellow}};
			\node[text=badark, font=\fontsize{6.5}{8}\selectfont]          at (7.84,0.64) {\texttt{\#FDE74C}};

			% Row labels
			\node[text=badark, font=\fontsize{8}{10}\selectfont\bfseries, anchor=east] at (-0.15,2.23) {Primär};
			\node[text=badark, font=\fontsize{8}{10}\selectfont\bfseries, anchor=east] at (-0.15,0.65) {Akzent};
		\end{tikzpicture}
	\end{center}

	\vspace{0.5em}
	\begin{columns}[T]
		\begin{column}{0.5\textwidth}
			\textbf{Schriftart:} Blinker (SIL Open Font License)\\[3pt]
			{\BlinkHeadline SemiBold · }{\BlinkRegular Regular · }\textit{Light (Kursiv)}
		\end{column}
		\begin{column}{0.5\textwidth}
			\textbf{Engine:} XeLaTeX mit \texttt{fontspec}\\[3pt]
			\textbf{Theme:} Beamer Metropolis · \texttt{progressbar=foot}
		\end{column}
	\end{columns}
\end{frame}

% ── Folie 3: Textbausteine & Typografie ──────────────────────────────────────
\begin{frame}{Textbausteine \& Typografie}

	\begin{columns}[T]
		\begin{column}{0.48\textwidth}

			\begin{block}{Normaler Block (\texttt{block})}
				Für neutrale Informationen und Fakten.\\[2pt]
				\textbf{Fett} · \textit{Kursiv} · \texttt{Monospace}
			\end{block}

			\vspace{0.4em}

			\begin{alertblock}{\textcolor{white}{\textbf{Alert-Block (\texttt{alertblock})}}}
				Für Fazit-Statements und Warnungen.
			\end{alertblock}

			\vspace{0.4em}

			\begin{exampleblock}{Beispiel-Block (\texttt{exampleblock})}
				Für Quellen, Hinweise, Ergänzungen.
			\end{exampleblock}

			\vspace{0.4em}
			\textbf{Farben:}
			\textcolor{bared}{Rot} ·
			\textcolor{baorange}{Orange} ·
			\textcolor{bablue}{Blau} ·
			\textcolor{bateal}{Teal}

		\end{column}
		\begin{column}{0.48\textwidth}

			\textbf{Aufzählungsliste:}
			\begin{itemize}
				\setlength{\itemsep}{0.35em}
				\item Erster Punkt mit \textbf{Hervorhebung}
				\item Zweiter Punkt — \textcolor{bared}{in Rot}
				\item Dritter Punkt — \textcolor{baorange}{in Orange}
				\item Vierter Punkt mit Quelle\footnotemark
			\end{itemize}

			\vspace{0.5em}

			\textbf{Fußnoten-Pattern} (overlay-sicher):
			\begin{beamercolorbox}[rounded=true,shadow=false,wd=\linewidth]{block body}
				\vspace{3pt}
				\fontsize{8}{10}\selectfont
				\texttt{\textbackslash footnotemark} im Text\\[2pt]
				\texttt{\textbackslash only<2->\{\textbackslash footnotetext[n]\{...\}\}}
				\vspace{3pt}
			\end{beamercolorbox}

			\vspace{0.5em}
			\textbf{Margins:}
			\texttt{\textbackslash setbeamersize\{text margin left=5pt\}}

			\footnotetext[1]{Quelle: Beispiel-Fußnote}

		\end{column}
	\end{columns}

\end{frame}

% ── Folie 4: Split-Slide — Text links, Bild rechts ──────────────────────────
\begin{frame}{\texttt{\textbackslash splitslide} — Text + Bild}

	\splitslide{%
		\begin{block}{Beschreibung}
			\vspace{0.2em}
			\texttt{\textbackslash splitslide\{text\}\{bild\}\{breite\}}
			\vspace{0.4em}

			Zweispaltiges Layout mit Text links und Bild rechts.
			Ideal für erklärende Folien mit visueller Unterstützung.
			\vspace{0.2em}
		\end{block}

		\vspace{0.5em}

		\begin{itemize}
			\setlength{\itemsep}{0.3em}
			\item Text und Bild gleichberechtigt
			\item Bild mit \texttt{roundedimage} (2:1)
			\item Breite relativ zur Spalte
		\end{itemize}
	}{pix/mining.jpg}{0.88}

\end{frame}

% ── Folie 5: Split-Slide mit Bild links ─────────────────────────────────────
\begin{frame}{\texttt{\textbackslash splitslideimgleft} — Bild links}

	\splitslideimgleft{pix/missverstaendnisse.jpg}{0.88}{%
		\begin{block}{Umgekehrtes Layout}
			\vspace{0.2em}
			Bild auf der linken Seite, Text rechts.
			\vspace{0.2em}
		\end{block}

		\vspace{0.5em}

		\texttt{\textbackslash splitslideimgleft\{bild\}\{breite\}\{text\}}

		\vspace{0.5em}

		\begin{itemize}
			\setlength{\itemsep}{0.3em}
			\item Variante von \texttt{\textbackslash splitslide}
			\item Für visuelle Abwechslung
			\item Bild-Argument kommt zuerst
		\end{itemize}
	}

\end{frame}

% ── Folie 6: Key-Point / Statement ──────────────────────────────────────────
\begin{frame}{\texttt{\textbackslash keypoint} — Zentrale Aussage}

	\keypoint{Bitcoin ist Freiheit.}

	\vspace{-1em}
	{\centering\fontsize{7.5}{10}\selectfont\color{bagray}
		\texttt{\textbackslash keypoint\{Bitcoin ist Freiheit.\}}\par}

\end{frame}

% ── Folie 7: Key-Point mit Farbe ────────────────────────────────────────────
\begin{frame}{\texttt{\textbackslash keypointcolor} — Akzentfarbe}

	\keypointcolor{baorange}{21 Millionen. Nicht mehr.}

	\vspace{-1em}
	{\centering\fontsize{7.5}{10}\selectfont\color{bagray}
		\texttt{\textbackslash keypointcolor\{baorange\}\{21 Millionen. Nicht mehr.\}}\par}

\end{frame}

% ── Folie 8: Vollbild mit Unterschrift ───────────────────────────────────────
\begin{frame}{\texttt{\textbackslash fullimage} — Bild mit Unterschrift}

	\fullimage{pix/mining.jpg}{Quelle: Beispielbild — Bitcoin-Mining-Anlage}

	\vspace{-0.5em}
	{\centering\fontsize{7.5}{10}\selectfont\color{bagray}
		\texttt{\textbackslash fullimage\{pix/bild.jpg\}\{Unterschrift\}}\par}

\end{frame}

% ── Folie 9: Vollbild ohne Unterschrift ──────────────────────────────────────
\begin{frame}{\texttt{\textbackslash fullimagenocaption} — Maximales Bild}

	\fullimagenocaption{pix/missverstaendnisse.jpg}

\end{frame}

% ── Folie 10: Zitat ─────────────────────────────────────────────────────────
\begin{frame}{\texttt{\textbackslash baquote} — Zitat}

	\baquote{The Times 03/Jan/2009 Chancellor on brink of second bailout for banks.}{Satoshi Nakamoto, Bitcoin Genesis Block}

	\vspace{-0.5em}
	{\centering\fontsize{7.5}{10}\selectfont\color{bagray}
		\texttt{\textbackslash baquote\{Zitat\}\{Quelle\}}\par}

\end{frame}

% ── Folie 11: Nummerierte Liste ──────────────────────────────────────────────
\begin{frame}{\texttt{\textbackslash numberedlist} — Nummerierte Liste}

	\twocol{%
		\numberedlist{Schritte zur Selbstverwahrung}{
			\item Bitcoin-Wallet installieren
			\item Seed-Phrase sicher aufbewahren
			\item Erste Transaktion empfangen
			\item Hardware-Wallet einrichten
		}
	}{%
		\vspace{0.5em}
		\texttt{\textbackslash numberedlist\{Titel\}\{\% \\
			\quad\textbackslash item Erster Schritt \\
			\quad\textbackslash item Zweiter Schritt \\
		\}}

		\vspace{1em}
		Erzeugt einen Block mit nummerierter Liste.
		Konsistentes Spacing via \texttt{itemsep}.
	}

\end{frame}

% ── Folie 12: Zwei-Spalten-Text ──────────────────────────────────────────────
\begin{frame}{\texttt{\textbackslash twocol} — Zwei Spalten}

	\twocol{%
		\begin{block}{Linke Spalte}
			\vspace{0.2em}
			Beliebiger Inhalt: Text, Listen, Blöcke, Bilder.
			\vspace{0.2em}
		\end{block}

		\vspace{0.4em}
		\texttt{\textbackslash twocol\{links\}\{rechts\}}
	}{%
		\begin{exampleblock}{Rechte Spalte}
			\vspace{0.2em}
			Einfacher als manuelles \texttt{columns}-Environment.
			Beide Spalten je 48\,\% Breite.
			\vspace{0.2em}
		\end{exampleblock}

		\vspace{0.4em}
		Ideal für Gegenüberstellungen und ergänzende Inhalte.
	}

\end{frame}

% ── Folie 13: Drei Spalten ───────────────────────────────────────────────────
\begin{frame}{\texttt{\textbackslash threecol} — Drei Spalten}

	\threecol{%
		\begin{block}{Spalte 1}
			\vspace{0.2em}
			\centering
			{\BlinkHeadline\fontsize{20}{24}\selectfont\color{bared}2009\par}
			\vspace{0.3em}
			Bitcoin-Whitepaper und Genesis Block
			\vspace{0.2em}
		\end{block}
	}{%
		\begin{block}{Spalte 2}
			\vspace{0.2em}
			\centering
			{\BlinkHeadline\fontsize{20}{24}\selectfont\color{baorange}21M\par}
			\vspace{0.3em}
			Maximale Anzahl an Bitcoin
			\vspace{0.2em}
		\end{block}
	}{%
		\begin{block}{Spalte 3}
			\vspace{0.2em}
			\centering
			{\BlinkHeadline\fontsize{20}{24}\selectfont\color{bablue}10\,min\par}
			\vspace{0.3em}
			Durchschnittliche Blockzeit
			\vspace{0.2em}
		\end{block}
	}

	\vspace{1em}
	{\centering\fontsize{7.5}{10}\selectfont\color{bagray}
		\texttt{\textbackslash threecol\{Spalte 1\}\{Spalte 2\}\{Spalte 3\}} — je 31\,\% Breite\par}

\end{frame}

% ── Folie 14: Vergleich / Comparison ─────────────────────────────────────────
\begin{frame}{\texttt{\textbackslash comparison} — Gegenüberstellung}

	\comparison{Fiat-Geld}{%
		\item Zentral kontrolliert
		\item Inflation durch Gelddrucken
		\item Kontrahentenrisiko
		\item Zensierbar
	}{Bitcoin}{%
		\item Dezentral und offen
		\item Feste Geldmenge (21M)
		\item Selbstverwahrung möglich
		\item Zensurresistent
	}

	\vspace{0.5em}
	{\centering\fontsize{7.5}{10}\selectfont\color{bagray}
		\texttt{\textbackslash comparison\{Titel L\}\{Items L\}\{Titel R\}\{Items R\}}\par}

\end{frame}

% ── Folie 15: Highlight-Box ─────────────────────────────────────────────────
\begin{frame}{\texttt{\textbackslash highlightbox} — Hervorhebung}

	\twocol{%
		\highlightbox{bared}{%
			\textbf{Wichtig:} Not your keys, not your coins!
			Verwahre deine Bitcoin immer selbst.
		}

		\vspace{0.6em}

		\highlightbox{baorange}{%
			\textbf{Tipp:} Hardware-Wallets bieten die beste
			Balance aus Sicherheit und Benutzerfreundlichkeit.
		}

		\vspace{0.6em}

		\highlightbox{bablue}{%
			\textbf{Info:} Bitcoin Austria bietet regelmäßige
			Workshops zur Selbstverwahrung an.
		}
	}{%
		\vspace{0.5em}
		\texttt{\textbackslash highlightbox\{farbe\}\{inhalt\}}

		\vspace{1em}
		Verfügbare Farben:
		\begin{itemize}
			\setlength{\itemsep}{0.3em}
			\item \textcolor{bared}{\textbf{bared}} — Warnungen
			\item \textcolor{baorange}{\textbf{baorange}} — Tipps
			\item \textcolor{bablue}{\textbf{bablue}} — Informationen
			\item \textcolor{bateal}{\textbf{bateal}} — Ergänzungen
		\end{itemize}
	}

\end{frame}

% ── Folie 16: Makro-Übersicht ────────────────────────────────────────────────
\begin{frame}{Makro-Übersicht}

	\vspace{0.3em}
	{\fontsize{8}{11}\selectfont
		\begin{tabular}{@{}l l@{}}
			\toprule
			\textbf{Makro} & \textbf{Verwendung} \\
			\midrule
			\texttt{\textbackslash splitslide\{text\}\{bild\}\{w\}}           & Text links, Bild rechts \\
			\texttt{\textbackslash splitslideimgleft\{bild\}\{w\}\{text\}}    & Bild links, Text rechts \\
			\texttt{\textbackslash keypoint\{Aussage\}}                        & Große zentrale Aussage \\
			\texttt{\textbackslash keypointcolor\{farbe\}\{Aussage\}}          & Aussage in Akzentfarbe \\
			\texttt{\textbackslash fullimage\{bild\}\{caption\}}               & Großes Bild + Unterschrift \\
			\texttt{\textbackslash fullimagenocaption\{bild\}}                 & Maximales Bild ohne Text \\
			\texttt{\textbackslash baquote\{Zitat\}\{Quelle\}}                & Zitatblock \\
			\texttt{\textbackslash numberedlist\{Titel\}\{items\}}             & Nummerierte Liste im Block \\
			\texttt{\textbackslash twocol\{links\}\{rechts\}}                  & Zwei-Spalten-Layout \\
			\texttt{\textbackslash threecol\{1\}\{2\}\{3\}}                    & Drei-Spalten-Layout \\
			\texttt{\textbackslash comparison\{T1\}\{I1\}\{T2\}\{I2\}}        & Gegenüberstellung \\
			\texttt{\textbackslash highlightbox\{farbe\}\{inhalt\}}            & Farbige Hervorhebung \\
			\texttt{\textbackslash roundedimage\{bild\}\{w\}\{ratio\}}         & Bild mit runden Ecken \\
			\texttt{\textbackslash bitcointitle\{Titel\}\{Untertitel\}}        & Titelfolie-Setup \\
			\texttt{\textbackslash bitcointitlegraphic\{bild\}}                & Titelbild \\
			\bottomrule
		\end{tabular}
	}

\end{frame}

\end{document}
